\subsection{Ordered Sets, Well-Orders, and Lattices}

\begin{definition} \label{def:ordered-set-vocabulary} \hlt{Ordered Set Vocabulary}\\
Let $S$ be a set, let $E,C\subseteq S$, and let $\leq_S\subseteq S^2$. Set $x<_Sy\Longleftrightarrow(x\leq_Sy\land x\neq y)$.\\
The relation $\leq_S$ is \hlt{antisymmetric} iff $(\forall x,y\in S)\,[(x\leq_Sy\land y\leq_Sx)\Rightarrow x=y]$. Elements $x,y\in S$ are \hlt{comparable} iff $x\leq_Sy\lor y\leq_Sx$.
\begin{enumerate}[label=(\roman*)]
\setlength{\itemsep}{0pt}
\item \hlt{Partial order}: The relation $\leq_S$ is a partial order iff it is reflexive, antisymmetric, and transitive. Then $(S,\leq_S)$ is a \hlt{poset}.
\item \hlt{Total order}: A partial order $\leq_S$ is total iff $(\forall x,y\in S)\,(x\leq_Sy\lor y\leq_Sx)$. Then $(S,\leq_S)$ is an \hlt{ordered set}.
\item \hlt{Chain}: If $\leq_S$ is a partial order, then $C$ is a chain iff $(\forall x,y\in C)\,(x\leq_Sy\lor y\leq_Sx)$.
\item \hlt{Well-order}: The order $\leq_S$ well-orders $S$ iff it is total and the following condition holds. $(\forall Y\in\mathcal{P}(S)\setminus\{\emptyset\})(\exists \ell\in Y)(\forall y\in Y)\,\ell\leq_Sy$.
\item \hlt{Bounds}: An element $\beta\in S$ is an upper bound of $E$ iff $(\forall x\in E)\,x\leq_S\beta$. It is a lower bound iff $(\forall x\in E)\,\beta\leq_Sx$.
\item \hlt{Boundedness}: The set $E$ is bounded above iff $(\exists \beta\in S)(\forall x\in E)\,x\leq_S\beta$. It is bounded below iff $(\exists \beta\in S)(\forall x\in E)\,\beta\leq_Sx$.
\item \hlt{Supremum}: For $\alpha\in S$, $\alpha=\sup E$ iff the following two conditions hold.
\begin{enumerate}[label=\arabic*.]
\setlength{\itemsep}{0pt}
\item \hlt{Upper bound}: $(\forall x\in E)\,x\leq_S\alpha$.
\item \hlt{Upper bound}: $(\forall x\in E)\,x\leq_S\alpha$.
\item \hlt{Least}: $(\forall \beta\in S)\,[((\forall x\in E)\,x\leq_S\beta)\Rightarrow\alpha\leq_S\beta]$.
\end{enumerate}
\item \hlt{Infimum}: For $\gamma\in S$, $\gamma=\inf E$ iff the following two conditions hold.
\begin{enumerate}[label=\arabic*.]
\setlength{\itemsep}{0pt}
\item \hlt{Lower bound}: $(\forall x\in E)\,\gamma\leq_Sx$.
\item \hlt{Greatest}: $(\forall \beta\in S)\,[((\forall x\in E)\,\beta\leq_Sx)\Rightarrow\beta\leq_S\gamma]$.
\end{enumerate}
\item \hlt{Maximal and minimal elements}: An element $m\in E$ is maximal iff $(\nexists x\in E)\,m<_Sx$. It is minimal iff $(\nexists x\in E)\,x<_Sm$.
\item \hlt{Least-upper-bound property}: This property is characterized by the following condition.\\
$(\forall A\in\mathcal{P}(S)\setminus\{\emptyset\})\,[((\exists \beta\in S)(\forall x\in A)\,x\leq_S\beta)\Rightarrow(\exists \alpha\in S)\,\alpha=\sup A]$.
\end{enumerate}
\end{definition}

\begin{illustrationfigure}[H]
  \centering
  \begin{tikzpicture}[
  x=0.80cm,
  y=0.82cm,
  >=Stealth,
  every node/.style={font=\small},
  axis/.style={gray!70, line width=0.55pt},
  setbar/.style={blue!70!black, line width=2.3pt, line cap=round},
  bound/.style={gray!50, line width=0.5pt},
  setpoint/.style={circle, fill=blue!70!black, inner sep=1.25pt},
  arrow/.style={->, gray!65, line width=0.5pt}
]
  \draw[axis] (-4.7,0) -- (4.7,0);

  \draw[setbar] (-1.35,0) -- (1.55,0);
  \foreach \x in {-0.75,-0.25,0.3,0.8,1.2} {
    \node[setpoint] at (\x,0) {};
  }

  \draw[arrow] (-1.35,0.60) -- (-1.35,0.16);
  \draw[arrow] (1.55,0.60) -- (1.55,0.16);
  \node[above] at (-1.35,0.60) {$\inf A$};
  \node[above] at (1.55,0.60) {$\sup A$};
  \node[below, yshift=-5pt] at (0.1,0) {$A$};

  \draw[bound] (-4.05,-0.42) -- (-4.05,-0.24) -- (-1.35,-0.24) -- (-1.35,-0.42);
  \node[below] at (-2.7,-0.42) {lower bounds};

  \draw[bound] (1.55,-0.42) -- (1.55,-0.24) -- (4.05,-0.24) -- (4.05,-0.42);
  \node[below] at (2.82,-0.42) {upper bounds};
\end{tikzpicture}

  \caption{The supremum is the least upper bound, while the infimum is the greatest lower bound.}
  \label{fig:supremum-infimum-bounds}
\end{illustrationfigure}

Figure \ref{fig:supremum-infimum-bounds} illustrates the bounds in Definition \ref{def:ordered-set-vocabulary} for a set $A$.\\
It marks $\inf A$ and $\sup A$ as extremal bounds; Example \ref{ex:supremum-need-not-be-attained} shows that neither need belong to the set.

\begin{example} \label{ex:supremum-need-not-be-attained} \hlt{Suprema Need Not Be Elements}\\
In $\Q$, set $E_1\defeq\{r\in\Q:r<0\}$ and $E_2\defeq\{r\in\Q:r\leq0\}$. Both sets have supremum $0$, but $0\notin E_1$ and $0\in E_2$.\\
For $E\defeq\{\frac{1}{n}:n\in\N\}$, $\sup E=1\in E$ and $\inf E=0\notin E$.\\
\hlt{Key Insight}: An extremal bound need not be an element of the set.
\end{example}

\begin{theorem} \label{thm:lub-property-gives-glb-property} \hlt{Least Upper Bounds Give Greatest Lower Bounds}\\
Let $(S,\leq)$ be an ordered set with the least-upper-bound property. Let $B\subseteq S$ be nonempty and bounded below. Then $\inf B$ exists in $S$.
\begin{proof} Set $L\defeq\{\ell\in S:(\forall b\in B)\,\ell\leq b\}$. $L\neq\emptyset$ by the boundedness of $B$ below. Choose $b_0\in B$.\\
Then $(\forall \ell\in L)\,\ell\leq b_0$, so $L$ is bounded above. Set $\alpha\defeq\sup L$; it exists in $S$ by the least-upper-bound property.\\
For $b\in B$, the element $b$ is an upper bound of $L$, so $\alpha\leq b$. Hence $\alpha\in L$. Moreover, $(\forall \ell\in L)\,\ell\leq\alpha$ by $\alpha=\sup L$.\\
Thus $\alpha$ is the greatest lower bound of $B$, and $\alpha=\inf B$.\\
\hlt{Structural Commentary}: The proof takes the supremum of the lower-bound set.\\
\hlt{Key Insight}: A nonempty set bounds its lower bounds above.
\end{proof}
\end{theorem}

\begin{definition} \label{def:lattice-vocabulary} \hlt{Lattice Vocabulary}\\
Let $(L,\leq)$ be a partially ordered set.
\begin{enumerate}[label=(\roman*)]
\setlength{\itemsep}{0pt}
\item \hlt{Meet}: For $x,y\in L$, the meet is $x\wedge y\defeq\inf\{x,y\}$ when this infimum exists.
\item \hlt{Join}: For $x,y\in L$, the join is $x\vee y\defeq\sup\{x,y\}$ when this supremum exists.
\item \hlt{Lattice}: The partially ordered set $L$ is a lattice iff the following condition holds. $(\forall x,y\in L)\,[x\wedge y\in L\land x\vee y\in L]$.
\item \hlt{Complete lattice}: A lattice $L$ is complete iff $(\forall E\subseteq L)\,[\inf E\in L\land\sup E\in L]$.
\item \hlt{Bottom}: An element $\bot\in L$ is a bottom element iff $(\forall x\in L)\,\bot\leq x$.
\item \hlt{Top}: An element $\top\in L$ is a top element iff $(\forall x\in L)\,x\leq\top$.
\item \hlt{Modular lattice}: A lattice $L$ is modular iff the following law holds. $(\forall x,y,z\in L)\,[x\leq z\Rightarrow z\wedge(y\vee x)=(z\wedge y)\vee x]$.
\end{enumerate}
\end{definition}
